\section{Elementary cases}\label{sec:first}

This section fixes the objects, derives the three counting identities that
every later argument uses, and settles the cases these identities settle
by themselves.  It then isolates, in the form of a lemma and a
counterexample, the exact point at which counting stops working; that
point is what the rest of the paper is about.  Two examples are introduced
here.

\subsection{The objects}

\begin{definition}\label{def:design}
A \emph{weighted design} of size $m$ and rank $k$ over $\FF$ is a pair
$\bigl((g_c)_{c\le m},(t_c)_{c\le m}\bigr)$ with $g_c\in\FF^k$, $t_c>0$
and $\sum_c t_c=1$, satisfying \eqref{eq:parseval}.  For
$C\subseteq\{1,\dots,m\}$ put $\SC:=\sum_{c\in C}g_c^{\vphantom{*}}g_c^{*}$;
the subset $C$ \emph{dominates} if $\SC-I_k\psd0$.  The assertion
$\Gtz(m,k)$ is that every design of size $m$ and rank $k$ over $\R$ admits
a dominating $k$-subset, and $\GtzC(m,k)$ is the same assertion over $\C$.
\end{definition}

Three features of the definition are deliberate.  The vectors are not
required to be distinct or nonzero, so repeated vectors and the zero
vector are admitted.  Excluding them would change the problem and would
not simplify it: a $k$-subset containing two parallel vectors is
rank-deficient and cannot dominate, and it is exactly this mechanism, seen
first in Example~\ref{ex:octa}, that makes selection combinatorial rather
than metric.  The weights are real even over $\C$.
And $\SC$ carries no weights: they enter only through
\eqref{eq:parseval}, which is why domination is preserved by any
operation that keeps the vectors and moves the weights.

Since $g_c^{\vphantom{*}}g_c^{*}\psd0$ for every $c$, adding indices to
$C$ can only increase $\SC$ in the Loewner order; a superset of a
dominating set dominates.  The cardinality restriction
$\lvert C\rvert=k$ is therefore the whole of the problem, and it is
sharp on both sides: fewer than $k$ vectors span a proper subspace and
leave $\SC$ singular, while at $\lvert C\rvert=m$ domination is automatic
(Lemma~\ref{lem:coparseval}).

\begin{lemma}[Lean]\label{lem:span}
The vectors of a design span $\FF^k$.  In particular $k\le m$, and there
is no design with $m<k$.
\end{lemma}

\begin{proof}
Let $u\in\FF^k$ satisfy $\langle u,g_c\rangle=0$ for every $c$.  Pairing
\eqref{eq:parseval} with $u$ gives
\[
  \lVert u\rVert^2\;=\;u^{*}I_ku
  \;=\;\sum_c t_c\,u^{*}g_c^{\vphantom{*}}g_c^{*}u
  \;=\;\sum_c t_c\,\lvert\langle u,g_c\rangle\rvert^2\;=\;0 ,
\]
so $u=0$.  The orthogonal complement of the span is trivial, and a
spanning family in $\FF^k$ has at least $k$ members.
\end{proof}

\subsection{Three counting identities}

Taking the trace of \eqref{eq:parseval} gives the identity on which every
counting argument in the subject rests.

\begin{lemma}[Lean]\label{lem:trace}
$\sum_c s_c=\sum_c t_c\ell_c=k$.  Hence $\max_c \ell_c\ge k$.
\end{lemma}

\begin{proof}
$\tr\bigl(g_c^{\vphantom{*}}g_c^{*}\bigr)=\lVert g_c\rVert^2=\ell_c$, and
$\tr I_k=k$; take traces in \eqref{eq:parseval} term by term.  The
weights are positive and sum to $1$, so $k=\sum_c t_c\ell_c$ is a
weighted mean of the $\ell_c$ and does not exceed their maximum.
\end{proof}

The leverage scores are also bounded above, individually.

\begin{lemma}[Lean]\label{lem:share}
$0\le s_c\le1$ for every $c$.  With Lemma~\ref{lem:trace} the leverage
scores therefore lie in $[0,1]$ and sum to $k$, which gives $m\ge k$ a
second time, with equality only if $s_c=1$ throughout.
\end{lemma}

\begin{proof}
Fix $e$ and separate the $e$-th term of \eqref{eq:parseval}:
\begin{equation}\label{eq:erased}
  I_k-t_e\,g_e^{\vphantom{*}}g_e^{*}
  \;=\;\sum_{c\ne e}t_c\,g_c^{\vphantom{*}}g_c^{*}\;\psd\;0 .
\end{equation}
Evaluating the quadratic form of \eqref{eq:erased} at $g_e$ gives
$0\le\ell_e-t_e\ell_e^2$.  If $\ell_e=0$ then $s_e=0\le1$; otherwise
divide by $\ell_e>0$.
\end{proof}

The third identity compares the unweighted sum over all indices with the
weighted one.

\begin{lemma}[Lean]\label{lem:coparseval}
For every design,
\begin{equation}\label{eq:coparseval}
  S_{\{1,\dots,m\}}-I_k\;=\;\sum_{c}(1-t_c)\,g_c^{\vphantom{*}}g_c^{*} .
\end{equation}
The right-hand side is positive semidefinite, and positive definite as
soon as $m\ge2$.
\end{lemma}

\begin{proof}
Subtract \eqref{eq:parseval} from $\sum_c g_c^{\vphantom{*}}g_c^{*}$.
Each $t_c$ is one summand of a sum of positive numbers equal to $1$, so
$t_c\le1$ and the coefficients $1-t_c$ are nonnegative; a nonnegative
combination of the positive semidefinite matrices
$g_c^{\vphantom{*}}g_c^{*}$ is positive semidefinite.  If $m\ge2$ then
$t_c<1$ for every $c$, since $t_c+t_{c'}\le1$ with $t_{c'}>0$ for any
second index $c'$.  Put $\sigma:=\min_c(1-t_c)/t_c>0$.  Every coefficient
of
\[
  \sum_c\bigl((1-t_c)-\sigma t_c\bigr)\,g_c^{\vphantom{*}}g_c^{*}
  \;=\;\sum_c(1-t_c)\,g_c^{\vphantom{*}}g_c^{*}\;-\;\sigma I_k
\]
is nonnegative by the choice of $\sigma$, the displayed equality being
\eqref{eq:parseval} scaled by $\sigma$.  Hence
$\sum_c(1-t_c)g_c^{\vphantom{*}}g_c^{*}\psd\sigma I_k\pd0$.
\end{proof}

The operator on the right of \eqref{eq:coparseval} is used again in
\S\ref{sec:chart}, where its inverse defines the pivot of a vector and
its positive definiteness supplies the whitening that makes duality
possible.  Nothing else in the present section needs it.

\subsection{The cases that close by counting}

\begin{proposition}[Lean]\label{prop:sizerank}
$\Gtz(k,k)$ and $\GtzC(k,k)$ hold.
\end{proposition}

\begin{proof}
By Lemma~\ref{lem:span} a design of rank $k$ has $m\ge k$ vectors, so
$m=k$ leaves exactly one $k$-subset, namely $\{1,\dots,m\}$, and it
dominates by Lemma~\ref{lem:coparseval}.
\end{proof}

\begin{proposition}[Lean]\label{prop:rankone}
$\Gtz(m,1)$ and $\GtzC(m,1)$ hold for every $m$.
\end{proposition}

\begin{proof}
At $k=1$ the matrices in \eqref{eq:parseval} are scalars and the
condition reads $\sum_c t_c\lvert g_c\rvert^2=1$; this is
Lemma~\ref{lem:trace} at $k=1$.  A weighted mean does not exceed its
maximum, so $\lvert g_{c_0}\rvert^2\ge1$ for some $c_0$, and the
singleton $\{c_0\}$, which has cardinality $k=1$, dominates.
\end{proof}

The two propositions dispose of the cells $m=k$ and $k=1$, and they are the
whole of what the counting identities give directly.  The range
\eqref{eq:elementary} in which the maximal-volume principle already
suffices is larger by exactly one column, $m=k+1$.  That column does reduce
to the scalar mechanism of Proposition~\ref{prop:rankone}, but only after a
change of variables which trades the vectors for the kernel of the
associated projection; it is Theorem~\ref{thm:corankone}, and it is where
the elementary methods end.

The classical statement is the uniform-weight instance, in the following
precise sense.

\begin{proposition}[Lean]\label{prop:classical}
Let $n\ge1$ and let $A\in\FF^{n\times k}$ satisfy $A^{*}A=I_k$.  If
$\Gtz(n,k)$ holds, respectively $\GtzC(n,k)$, then there is an injection
$\iota:\{1,\dots,k\}\to\{1,\dots,n\}$ with
\[
  (A_\iota)^{*}A_\iota^{\vphantom{*}}\;\psd\;\tfrac1n I_k,
  \qquad A_\iota:=\bigl(A_{\iota(1)},\dots,A_{\iota(k)}\bigr)^{*}\!,
\]
where $A_r^{*}$ denotes the $r$-th row of $A$.
\end{proposition}

\begin{proof}
Put $g_r:=\sqrt n\,A_r$ and $t_r:=1/n$.  Then
$\sum_r t_r g_r^{\vphantom{*}}g_r^{*}=\sum_r A_r^{\vphantom{*}}A_r^{*}
=A^{*}A=I_k$, so this is a design of size $n$ and rank $k$.  Let $C$
dominate and let $\iota$ enumerate $C$ in increasing order.  Then
$(A_\iota)^{*}A_\iota^{\vphantom{*}}=\sum_{r\in C}A_r^{\vphantom{*}}A_r^{*}
=\frac1n\SC$, so
\[
  (A_\iota)^{*}A_\iota^{\vphantom{*}}-\tfrac1n I_k
  \;=\;\tfrac1n\bigl(\SC-I_k\bigr)\;\psd\;0 . \qedhere
\]
\end{proof}

At a fixed size the implication goes only this way, so the weighted
statement is formally the stronger of the two; over all sizes they are
equivalent, for the reason given in \S\ref{sec:intro}.

\subsection{Where counting stops}\label{sec:stops}

Domination admits a formulation with no matrices in it.  A $k$-subset $C$
dominates if and only if
\begin{equation}\label{eq:pointwise}
  \sum_{c\in C}\lvert\langle x,g_c\rangle\rvert^2\;\ge\;\lVert x\rVert^2
  \qquad\text{for every }x\in\FF^k ,
\end{equation}
since the left-hand side is the quadratic form of $\SC$ at $x$.  In the
same language \eqref{eq:parseval} says that
$\sum_c t_c\lvert\langle x,g_c\rangle\rvert^2=\lVert x\rVert^2$ for every
$x$: the squared projections of a fixed direction on the vectors of a
design average to the squared length of that direction.  The averaging
mechanism of Proposition~\ref{prop:rankone} therefore survives in full
generality, one direction at a time.

\begin{lemma}[Lean]\label{lem:direction}
For every design and every $x\in\FF^k$ there is an index $c$ with
$\lvert\langle x,g_c\rangle\rvert^2\ge\lVert x\rVert^2$.
\end{lemma}

\begin{proof}
The numbers $\lvert\langle x,g_c\rangle\rvert^2$ have weighted mean
$\lVert x\rVert^2$ by \eqref{eq:parseval}, and a weighted mean with
positive weights summing to $1$ does not exceed the maximum of the
numbers averaged.
\end{proof}

At $k=1$ this is the theorem, because there is only one direction up to
scalars and the witnessing index may be chosen once and for all.  At
$k\ge2$ the witness of Lemma~\ref{lem:direction} depends on $x$, and
\eqref{eq:pointwise} demands a set of $k$ indices that works
simultaneously for every $x$ on a sphere of dimension $k-1$.  No
averaging identity produces such a set: the quantity to be dominated in
\eqref{eq:dominate} is a matrix, and while \eqref{eq:parseval} controls
every scalar linear functional of the family, the least eigenvalue is not
one of them.

The natural scalar relaxation makes the failure quantitative.  If $C$
dominates then all $k$ eigenvalues of $\SC$ are at least $1$, so
\begin{equation}\label{eq:tracetest}
  \tr\SC\;=\;\sum_{c\in C}\ell_c\;\ge\;k .
\end{equation}
This necessary condition is satisfied in every design by a subset that is
found without any work: by Lemma~\ref{lem:trace} some $\ell_{c_0}\ge k$,
and any $k$-subset containing $c_0$ satisfies \eqref{eq:tracetest}
already.  The trace test thus never fails to be satisfiable and never
distinguishes a dominating subset from a non-dominating one.  The
following design, which will also serve as a first illustration of the
cell $(6,3)$, shows how far apart the two can be.

\begin{example}\label{ex:octa}
Let $k=3$, let the six vectors be $\pm\sqrt3\,e_1$, $\pm\sqrt3\,e_2$,
$\pm\sqrt3\,e_3$ and let every weight be $\tfrac16$.  Then
\[
  \sum_c t_c\,g_c^{\vphantom{\mathsf{T}}}g_c\tp
  \;=\;\tfrac16\cdot3\cdot2\sum_{i=1}^{3}e_i^{\vphantom{\mathsf{T}}}e_i\tp\;=\;I_3 ,
\]
so this is a design at $(6,3)$, with $\ell_c=3$ and $s_c=\tfrac12$ for
every $c$, in accordance with Lemmas~\ref{lem:trace} and~\ref{lem:share}.
Every triple satisfies $\tr\SC=9\ge3$.  The triple
$\{\sqrt3e_1,-\sqrt3e_1,\sqrt3e_2\}$ nevertheless has
$\SC=6\,e_1^{\vphantom{\mathsf{T}}}e_1\tp+3\,e_2^{\vphantom{\mathsf{T}}}e_2\tp$, which is
singular in the direction $e_3$ and does not dominate; the triple
$\{\sqrt3e_1,\sqrt3e_2,\sqrt3e_3\}$ has $\SC=3I_3\pd I_3$ and dominates
strictly.  The trace is blind to the difference, and the difference is
which subspaces the chosen vectors miss.
\end{example}

Selection is therefore a combinatorial problem about which directions a
$k$-subset leaves uncovered, constrained by the metric identity
\eqref{eq:parseval}.  The remaining sections develop the two instruments
that address it: a change of coordinates that turns \eqref{eq:dominate}
into a statement about principal blocks of a projection
(\S\ref{sec:chart}), and, when that does not decide the cell, a
variational analysis of the worst design (\S\ref{sec:compactness} and
\S\ref{sec:interior}).

\subsection{The tetrahedron}

In Example~\ref{ex:octa} one triple dominates strictly.  The next design
admits no strict domination at all, and its equality is exact at every
triple.

\begin{example}[Lean]\label{ex:tetra}
Let $k=3$ and take the four vectors
\[
  g_0=(1,1,1),\quad g_1=(1,-1,-1),\quad g_2=(-1,1,-1),\quad g_3=(-1,-1,1)
\]
with equal weights $t_c=\tfrac14$.  Each diagonal entry of
$\sum_c g_c^{\vphantom{\mathsf{T}}}g_c\tp$ is $4$, and each off-diagonal entry is
$1-1-1+1=0$, so
\begin{equation}\label{eq:tetraparseval}
  \sum_{c=0}^{3} g_c^{\vphantom{\mathsf{T}}}g_c\tp\;=\;4I_3,
  \qquad\text{hence}\qquad
  \sum_{c=0}^{3}\tfrac14\,g_c^{\vphantom{\mathsf{T}}}g_c\tp\;=\;I_3 :
\end{equation}
a design at $(4,3)$.  Here $\ell_c=3$ and $s_c=\tfrac34$ for every $c$, so
Lemma~\ref{lem:trace} holds with equal terms, $4\cdot\tfrac34=3=k$, and
$\langle g_i,g_j\rangle=-1$ for $i\ne j$.  The four vectors are the unit
vectors to the vertices of a regular tetrahedron, scaled by $\sqrt3$.

Let $C=\{0,1,2,3\}\setminus\{c_0\}$ be one of the four triples.
Subtracting the omitted term from \eqref{eq:tetraparseval},
\[
  \SC\;=\;4I_3-g_{c_0}^{\vphantom{\mathsf{T}}}g_{c_0}\tp .
\]
The matrix $g_{c_0}^{\vphantom{\mathsf{T}}}g_{c_0}\tp$ has eigenvalue
$\ell_{c_0}=3$ on the line $\R g_{c_0}$ and $0$ on its orthogonal
complement, so
\begin{equation}\label{eq:tetraspec}
  \operatorname{spec}\SC=\{1,4,4\},\qquad \lmin(\SC)=1 .
\end{equation}
Every triple dominates, and every triple dominates with least eigenvalue
exactly $1$.  For $c_0=0$ the certificate is a sum of squares: writing
$x=(x_1,x_2,x_3)$,
\[
  x\tp\bigl(\SC-I_3\bigr)x
  \;=\;3\lVert x\rVert^2-\langle x,g_0\rangle^2
  \;=\;3\sum_i x_i^2-\Bigl(\sum_i x_i\Bigr)^{\!2}
  \;=\;\sum_{i<j}(x_i-x_j)^2 ,
\]
which vanishes exactly on $\R g_0$, the direction omitted from $C$.
\end{example}

The tetrahedron is the smallest design at rank three in which
\eqref{eq:dominate} holds with equality.  It is the equality case
of the corank-one classification: the tie law \eqref{eq:tielaw} of
\S\ref{sec:chart} reads $3\cdot\tfrac14\cdot3=\tfrac94=3-1+\tfrac14$
here.  Splitting its vectors produces extremal designs at every size
$m\ge4$ and rank $3$, which is how the no-margin principle of
\S\ref{sec:ties} reaches the open cells.  Its volume-sampling mixture is
computed in \S\ref{sec:ties} and equals $(x-1)(x-4)^2$, with least root
exactly at the threshold.  And it is the calibration point of the
first-order system of \S\ref{sec:interior}, where the multiplier is
$\Lambda=\tfrac13I_3$ with $v=1$ and each of the four triples carries
multiplier $\tfrac14$.

\begin{figure}[ht]
  \centering
  % The regular tetrahedron design at (4,3), and the spectrum of one of its
% triples.  Orthographic projection of R^3 with azimuth 15 and elevation 70
% degrees, scaled by 1.35cm:
%   X = 1.35(-0.2588 x_1 + 0.9659 x_2)
%   Y = 1.35(-0.9077 x_1 - 0.2432 x_2 + 0.3420 x_3)
\begin{tikzpicture}[
  x=1cm, y=1cm,
  vec/.style={-{Latex[length=2mm,width=1.5mm]}, line width=0.75pt},
  cube/.style={densely dotted, line width=0.3pt, black!35},
  edge/.style={line width=0.45pt, black!65},
  hidden/.style={line width=0.45pt, black!45, dashed, dash pattern=on 1.4pt off 1.2pt},
  every node/.style={font=\small}]

% ---- the eight vertices of [-1,1]^3 -------------------------------------
\coordinate (ppp) at ( 0.9546,-1.0920);   % g_0 = ( 1, 1, 1)
\coordinate (ppm) at ( 0.9546,-2.0154);   %       ( 1, 1,-1)
\coordinate (pmp) at (-1.6534,-0.4354);   %       ( 1,-1, 1)
\coordinate (pmm) at (-1.6534,-1.3588);   % g_1 = ( 1,-1,-1)
\coordinate (mpp) at ( 1.6534, 1.3588);   %       (-1, 1, 1)
\coordinate (mpm) at ( 1.6534, 0.4354);   % g_2 = (-1, 1,-1)
\coordinate (mmp) at (-0.9546, 2.0154);   % g_3 = (-1,-1, 1)
\coordinate (mmm) at (-0.9546, 1.0920);   %       (-1,-1,-1)
\coordinate (O)   at ( 0,0);

% ---- the cube, as a frame of reference ----------------------------------
\draw[cube] (ppp)--(ppm) (ppp)--(pmp) (ppp)--(mpp)
            (ppm)--(pmm) (ppm)--(mpm) (pmp)--(pmm) (pmp)--(mmp)
            (mpp)--(mpm) (mpp)--(mmp) (pmm)--(mmm) (mpm)--(mmm) (mmp)--(mmm);

% ---- the tetrahedron on the four even-sign vertices ----------------------
\draw[hidden] (pmm)--(mpm);
\draw[edge]   (ppp)--(pmm) (ppp)--(mpm) (mpm)--(mmp) (mmp)--(pmm) (ppp)--(mmp);

% ---- the four vectors of the design -------------------------------------
\draw[vec] (O)--(ppp) node[below right=-1pt] {$g_0$};
\draw[vec] (O)--(pmm) node[below left=-1pt]  {$g_1$};
\draw[vec] (O)--(mpm) node[right=1pt]        {$g_2$};
\draw[vec] (O)--(mmp) node[above left=-1pt]  {$g_3$};
\fill (O) circle (1.1pt);

% ---- the spectrum of a triple -------------------------------------------
% the eigenvalue lambda is drawn at height  -1.6 + 0.64 lambda
\begin{scope}[shift={(4.55,0)}]
  \draw[line width=0.5pt] (0,-1.66)--(0,1.72);
  \foreach \l/\lab in {0/0,1/1,2/2,3/3,4/4,5/5}{
    \pgfmathsetmacro{\hh}{-1.6+0.64*\l}
    \draw[line width=0.5pt] (-0.07,\hh)--(0.07,\hh);
    \node[anchor=east, font=\footnotesize] at (-0.11,\hh) {\lab};
  }
  % the threshold: the eigenvalue of I_3
  \draw[dashed, dash pattern=on 1.7pt off 1.5pt, line width=0.5pt]
        (0.06,-0.96)--(1.05,-0.96);
  \node[anchor=west, font=\footnotesize] at (1.08,-0.96) {$\lmin(\SC)=1$};
  % the three eigenvalues 1, 4, 4
  \fill (0.30,-0.96) circle (1.7pt);
  \fill (0.30, 0.96) circle (1.7pt);
  \fill (0.62, 0.96) circle (1.7pt);
  \node[anchor=west, font=\footnotesize] at (0.78,0.96) {twice};
  \node[anchor=south] at (0.45,1.80) {$\operatorname{spec}\SC$};
\end{scope}

\end{tikzpicture}

  \caption{The design of Example~\ref{ex:tetra}: four vectors to the
  vertices of a regular tetrahedron, drawn inside the cube whose even-sign
  vertices they are.  Every triple omits one vector and has spectrum
  $\{1,4,4\}$ by \eqref{eq:tetraspec}, so the least eigenvalue sits
  exactly on the threshold.}
  \label{fig:tetrahedron}
\end{figure}
